\section{Conclusion}
\label{sec:conclusion}

Lead-lag analysis in iterative training collapse is far more fragile than the current literature assumes.
We demonstrated three pitfalls that undermine standard approaches: method choice alone swings FDR significance from 100\% to 25\% ($\kappa{=}0.185$), first-differencing exposes canary$\to$perplexity associations as entirely spurious ($F{\approx}0$, $p{=}0.99$), and regime dependence causes the same Granger test to shift from 24/24 to 0/8 significant pairs across contamination levels.
These findings establish that single-method, single-regime analyses---the current default---cannot reliably distinguish genuine temporal precedence from statistical artifacts.

Our multi-method consensus framework addresses this fragility by requiring convergent evidence across ${\geq}4/5$ independent methods and ${\geq}3/4$ random seeds.
Calibrated on built-in experimental controls, the framework achieves 0\% false positive rate at the null control ($\alpha{=}0.0$) and 0\% false negative rate against ground-truth pairs at full contamination ($\alpha{=}1.0$), with four tiers of practical guidance scaling from exploratory screening to publication-grade claims.

Applying the framework yields a concrete finding: token entropy serves as a one-generation-ahead canary for diversity collapse at $\alpha{=}1.0$ (with preliminary single-seed evidence at $\alpha{=}0.75$), with $|r|{>}0.99$ and complete directional asymmetry (0/24 reverse-direction FDR-significant pairs).
This result is deliberately scoped to lexical diversity, not model collapse broadly, since the perplexity channel proved entirely spurious under differencing.

The methodological pitfalls we expose are not specific to our experimental setting---they arise from generic properties of short, non-stationary time series---and the consensus framework transfers directly to any domain where multiple detection methods can be applied in parallel.
Limitations and remaining gaps are discussed in \S\ref{sec:discussion}.